\section{Accuracy and the sequence of seed passes}\label{sec:accuracy}
We first sum the node estimates to prove the refinement guarantee. We then
use that guarantee to improve the simultaneous seed factor and obtain the
answer to the original input. Throughout the first step the entire seed
table is fixed and satisfies the hypotheses of
Theorem~\ref{thm:refinement}.

\subsection{Summing the comparison regrets}
\begin{proof}[Proof of Theorem~\ref{thm:refinement}, accuracy part]
Fix a starting state with true cost $D>0$ and its comparison tree from
Section~\ref{sec:comparison}. For brevity write
$\overline d_t=\max(c,d_t)$ at every state, as in
Section~\ref{sec:regret}.  Terminal values are exact, so summing
\eqref{eq:node-regret-definition} over the tree at a fixed index gives
the identity
\[
 \overline d_t(\mathrm{root})
 =C_{\mathrm{term}}+L\sum_{\mathrm{nonterminal}}e
                      +\sum_{\mathrm{nonterminal}}\mathcal R_t.
\]
Average this identity over $t=1,\ldots,S$.  Since the root estimates
decrease with $t$, the final positive error satisfies the deterministic
bound
\begin{equation}\label{eq:root-regret-telescope}
 (d_S(\mathrm{root})-D)_+
 \le C_{\mathrm{term}}-D+L\sum_{\mathrm{nonterminal}}e
       +\frac1S\sum_{\mathrm{nonterminal}}\sum_{t=1}^S\mathcal R_t.
\end{equation}
Lemma~\ref{lem:comparison-tree} bounds the first two terms by
$(8MJ/P+4LJ/P)D$.  We bound the last sum using the three classes
of comparison nodes.

On bad nodes use Lemma~\ref{lem:generic-regret} and $W\le4Fc$.
Corollary~\ref{cor:bad-node-cost} bounds their total contribution by
\[
 40000F^2 H^{1/4}(1+J/P+J/T)D.
\]
On nodes that are not bad but satisfy $|\widehat v|>1-\tau$, the same
lemma gives $(H^{-1/4}+4\tau F)c$. Every other nonterminal satisfies
Lemma~\ref{lem:comparison-band}, so Proposition~\ref{prop:node-regret}
applies. Its dimension and mass bounds are
\[
 \ln d_{\mathcal C}\le4(p_*+21)\ell,\qquad h_\Sigma\le64Fb\le64Fc.
\]
The cost sum over all nonterminal levels is at most $2JD$.

The conditional estimate at a node is valid after fixing its complete
child tables. Taking expectations and then summing uses only the tower
property and linearity of expectation. Independence between comparison
nodes, or simultaneous success of their group samples, is unnecessary.
Combining these bounds, the expected excess at the starting state divided
by $D$ is at most
\begin{align}
 \mathcal E={}&(8M+512F)J/P
 +40000F^2H^{1/4}(1+J/P+J/T)/S \notag\\
 &+\frac{2J}{S}\left[
 H^{-1/4}+4\tau F
 +64F\left(\frac{4(p_*+21)\ell}{\eta}+2\kappa\right)
 +2\delta\right]
 +132FJH^{-3}.\label{eq:total-error}
\end{align}
Here the last term includes both numerical-mixture loss and group-failure
loss: they are at most $128FJH^{-3}$ and $4FJH^{-4}$, respectively.

We verify the numerical bound explicitly. Substituting
$J\le100H/\ell$, $T\ge H^{3/4}$, $S=H/\ell^2$, $F=\ell$,
$M\le H^{1/20}$ and $P\ge H^{20}$ yields
\begin{align*}
 \mathcal E\le{}&52000H^{-18}
 +8.04\cdot10^6\ell^4H^{-1/2}
 +200\ell H^{-1/4}\\
 &+\frac{10^6(p_*+21)}{\ell^2}+13200H^{-2}.
\end{align*}
For the fourth term we have included
$[26400+51200(p_*+21)]/\ell^2+400/\ell^9$ from the bracket in
\eqref{eq:total-error}. Each of the five displayed terms is at most
$1/(8\ell)$ at $\ell\ge2^{1000}$ and $H\ge2^{\ell/2}$.
For example, the second requires only
\[
 \log_2(8\cdot8.04\cdot10^6)+5\log_2\ell\le\ell/4;
\]
it holds at the cutoff and remains true thereafter. The other
negative-power terms require weaker inequalities of the same form.
For the fourth, use $p_*+21\le\sqrt\ell$ and
$\sqrt\ell\ge2^{500}>8\cdot10^6$.
Consequently $\mathcal E<1/\ell$.

Since $\overline d_S-D=(d_S-D)_+$, this is the required expected
positive-excess bound. The deterministic lower bound is the invariant
from Section~\ref{sec:bands}. When $D=0$, that invariant and
$w_0\le2Fc$ give $d_S=0$ directly, without dividing by $D$.
The numerical construction in Section~\ref{sec:numerics} supplies the
mixtures used in this proof, and Section~\ref{sec:work} verifies finite
execution of every defined operation.
\end{proof}

\subsection{Improving simultaneous seeds}
The initial table of Theorem~\ref{thm:seed} has factor $A_{\rm init}$.
The algorithm maintains a power-of-two factor $A$ and an integer seed
$U\in[0,N]$ on the full state domain. A seed is called \emph{good} when
$c\le U\le Ac$ at every state. This is a mathematical event, not a test
performed by the algorithm.

While $A>F$, set $a=A/F$ and use $H^2$ independent copies of the
refinement, conditional on the previous complete table. For each state,
form the $H^2$ integers $\lceil a d_S\rceil$, take their upper middle
order statistic, and clamp the result to $[0,N]$. This defines the next
seed. Update
\begin{equation}\label{eq:seed-factor-update}
 A\longleftarrow4a.
\end{equation}
All copies share the preceding table, but use independent new streams.
Entries of each copy and of the new seed are stored on demand; their
coherent full-table interpretation is justified in Section~\ref{sec:work}.

\begin{lemma}\label{lem:seed-amplification}
Conditional on any fixed good preceding seed, the next seed is good for
factor $4a$ except with probability at most $\exp(10H-H^2)$.
\end{lemma}
\begin{proof}
At a fixed state of positive integer cost $c$, the refinement theorem and
Markov's inequality give
\[
 \Pr[d_S>2c]\le1/\ell.
\]
Every copy has $\lceil a d_S\rceil\ge c$. A copy with $d_S\le2c$ has
$\lceil a d_S\rceil\le2ac$, because $a$ and $c$ are integers. If strictly
more than half the copies have this upper bound, their upper middle order
statistic does too, hence in particular is at most $4ac$. The probability
that at least half fail is at most
\[
 2^{H^2}\ell^{-H^2/2}\le e^{-H^2}.
\]
Clamping cannot spoil the lower bound $c\le N$ or the upper bound.
At zero-cost states every copy is zero under the good seed.

There are at most $2M\max(1,|x|)$ labeled nodes and at most
$(|y|+1)^2$ valid states per node. Their product is less than $e^{10H}$
in the large branch. Union over this full domain proves the assertion.
No independence between different states is used.
\end{proof}

The update preserves powers of two and strictly decreases $A$, since
$F>4$. Writing $s_{\rm pass}$ for the number of median passes plus the
single final pass below, we have
\begin{equation}\label{eq:pass-count}
 s_{\rm pass}
 =O\left(\frac{(p_*+1)\ell}{\log_2\ell}\right)<H.
\end{equation}
Indeed each median pass reduces $\log_2 A$ by $\log_2 F-2$, starting
from $(p_*+4)\log_2 H$ and stopping by $A\le F$.
Apply Lemma~\ref{lem:seed-amplification} to the first failed transition
after a good history. The total probability that any seed needed by the
algorithm is not good is less than
\begin{equation}\label{eq:seed-total-failure}
 0.01+H\exp(10H-H^2)<1/10.
\end{equation}
The initial $0.01$ is the simultaneous initial-table failure probability.
This estimate covers adaptive requests because every seed guarantee is
over its entire potential domain.

\subsection{The final answer}
When $A\le F$, use one new refinement with $a=1$ at the full root state
$(I,q)=((0,|x|],(0,|y|))$. Clamp its value to $[0,N]$ and round down to
obtain $\widehat D$. Fix any complete good seed before exposing this
final refinement's independent randomness. For the positive cost
$D=\ED(x,y)$ the refinement theorem gives the deterministic lower bound
$d_S\ge D$ under this conditioning and
\[
 \Pr[d_S>(1+\varepsilon)D\mid U\text{ fixed and good}]
 \le\frac{1}{\ell\varepsilon}\le\frac1{10}.
\]
Together with~\eqref{eq:seed-total-failure}, the unconditional failure
probability is less than $1/5$, and thus less than $1/3$. The lower bound
is not asserted on executions with a bad seed.

Clamping preserves the lower bound because $D\le N$. Rounding down also
preserves it because $D$ is an integer. Neither operation worsens the
upper bound. The initial equality test returns zero deterministically
when $D=0$, and the exact branches of Section~\ref{sec:geometry} are also
correct. This establishes the accuracy assertion of
Theorem~\ref{thm:main} on every input. It remains to prove that the
finite computations just specified, with all randomness and arithmetic
charged, have the claimed expected running time.
